\section{Estimating accuracy from few labels}
\label{sec:importance}

Labels are often the bottleneck for evaluating a deployed classifier. With class probabilities
the predicted uncertainty $u_i = 1 - p_{\max,i}$ can guide which items to label. We compare
uniform sampling with the confidence-guided proposal
$q_i = (1 - \varepsilon)\,u_i^{\alpha} / \sum_j u_j^{\alpha} + \varepsilon / N$ with
$\varepsilon = \ISEpsilon$ and $\alpha = \ISAlpha$~\citep{sawade2010are,kossen2021active},
estimating the error rate with the Hansen--Hurwitz estimator~\citep{hansen1943}, over
\ISTrials{} seeded trials per label budget; the sampler never sees a label it did not draw.

\begin{table}[t]
  \caption{RMSE of the accuracy estimate from $n$ labels, for uniform and confidence-guided
    sampling, and the ratio of their variances (uniform over guided; above 1 favours guided
    sampling), over \ISTrials{} trials.}
  \label{tab:importance}
  \centering\small
  \setlength{\tabcolsep}{4pt}
  \fittable{\begin{tabular}{lrrccccccccc}
\toprule
 & & & \multicolumn{3}{c}{$n = 25$} & \multicolumn{3}{c}{$n = 100$} & \multicolumn{3}{c}{$n = 400$} \\
\cmidrule(lr){4-6}\cmidrule(lr){7-9}\cmidrule(lr){10-12}
Dataset & Pool & Accuracy & uniform & guided & ratio & uniform & guided & ratio & uniform & guided & ratio \\
\midrule
German Credit & 1{,}000 & 0.471 & 0.097 & 0.212 & 0.21 & 0.049 & 0.105 & 0.22 & 0.025 & 0.052 & 0.23 \\
Breast Cancer Wisconsin & 569 & 0.427 & 0.101 & 0.090 & 1.27 & 0.048 & 0.045 & 1.14 & 0.024 & 0.021 & 1.32 \\
SST-2 & 872 & 0.960 & 0.039 & 0.057 & 0.46 & 0.020 & 0.029 & 0.47 & 0.010 & 0.014 & 0.48 \\
AG News & 7{,}600 & 0.879 & 0.067 & 0.135 & 0.24 & 0.032 & 0.069 & 0.22 & 0.017 & 0.032 & 0.28 \\
\bottomrule
\end{tabular}
}
\end{table}

The estimator is unbiased in practice (largest absolute bias \IsMaxAbsBias), but the guided
proposal reduces variance on only \IsDatasetsGuidedWins{} of the \NDatasets{} datasets
(Table~\ref{tab:importance}). At
\ISMaxBudget{} labels the ratio of uniform to guided variance is \IsVarRatioWdbcNFourZeroZero{}
on Breast Cancer Wisconsin and \IsVarRatioGermanNFourZeroZero, \IsVarRatioSstNFourZeroZero{} and
\IsVarRatioAgnewsNFourZeroZero{} on German Credit, SST-2 and AG News. The distribution itself
shows why. Jev rounds probabilities to two decimals, so \IsZeroUncertaintyAgnews{} of
\IsPoolAgnews{} AG News items have a reported uncertainty of exactly zero; they receive only the
minimum sampling probability yet contain \IsErrorsZeroUncertaintyAgnews{} of the
\IsErrorsAgnews{} errors, each weighted by $1/\varepsilon$. And for a 0/1 loss a proposal
proportional to the expected loss is not variance-optimal even under perfect calibration. The
check that reveals this needs only a small labelled sample (Appendix~\ref{app:importance},
Figure~\ref{fig:is}).
